\section{Лемма о почти перпендикуляре.}
\begin{lemma}[Лемма Рисса о почти перпендикуляре]
    Пусть $X$ -- нормированное пространство, $L \subset X$ -- его собственное замнкнутое подпространство.
    Тогда $\forall \epsilon > 0 \exists x_\epsilon \in X\colon \|x_\epsilon\| = 1$ и $\rho(x_\epsilon, L) \geq 1 -\epsilon$.
\end{lemma}
\begin{proof}
    Поскольку $L$ -- собственное подпространство, существует $z \in X \setminus L$.
    В силу замкнутости $\rho(z, L) = \inf_{x \in L} \|z - x\| > 0$.
    Заметим, что $\rho(z, L) \leq \|z\|$.
    Существует последовательность $\{y_n\} \subset L$ такая, что
    \[
        \rho(z, L) \leq \|z - y_n\| \leq \rho(z, L) + \frac{1}{n},
    \]
    то есть $\|z - y_n\| \rightarrow 0$.

    Тогда $\forall \epsilon > 0$ верно, что существует некоторый $n_\epsilon \in \N$ такой, что
    \[
        \frac{\rho(z, L)}{1 - \epsilon} \geq \frac{1}{n_\epsilon} + \rho(z, L) \geq \|z - y_{n_\epsilon}\| \geq \rho(z, L).
    \]
    Положим тогда $x_\epsilon = \frac{z - y_{n_\epsilon}}{\|z - y_{n_\epsilon}\|}$.
    Он имеет единичную норму и
    \[
        \rho(x_\epsilon, L) = \rho\left(\frac{z - y_{n_\epsilon}}{\|z - y_{n_\epsilon}\|}\right) = \frac{1}{\|z - y_{n_\epsilon}\|}\rho(z, L) \geq \frac{1 - \epsilon}{\rho(z, L)}\rho(z, L) = 1 - \epsilon.
    \]
    Что и требовалось в утверждении теоремы.
\end{proof}
\begin{theorem}[Теорема Рисса о некомпактности единичной сферы]
    Если $(X, \|\cdot\|)$ -- бесконечномерное нормированное пространство, то единичная сфера $S$ в нём не является компактной.
\end{theorem}
\begin{proof}
    Поскольку $X$ -- бесконечномерно, то существует некоторый вектор из единичной сферы $x_0$.
    Пусть $L_0 = \Lin\{x_0\}$.
    По теореме о почти перпендикуляре ($L_0$ -- замкнуто, поскольку конечномерно, и собственно по аналогичным причинам) существует $x_1\colon \rho(x_1, L_0) \geq \frac{1}{2}$ и $x_1 \in S$.
    Аналогично мы можем построить $x_2, x_3, \ldots$.
    Таким образом мы получили $\frac{1}{2}$-разреженную последовательность точек на единичной сфере.
    Но, если сфера -- компактна, то из этой последовательности мы смогли бы выделить сходящуюся подпоследовательность -- чего мы сделать, очевидно, не можем.
    Следовательно $S$ не является компактом.
\end{proof}
\section{Линейные функционалы}
\begin{definition}
    Пусть $X$ -- нормированное пространство.
    Линейное отображение
    \[
        f\colon X \to \R
    \]
    называется линейным функционалом.
\end{definition}
\begin{definition}
    Ядром линейного функционала $f$ называется $\ker f = f^{-1}(0)$.
\end{definition}
\begin{note}
    Пусть $f$ -- линейный функционал -- не тождественно равен 0.
    Тогда $\exists x_0 \in X \setminus \ker f$.

    Для всякого $x \in X$ 
    \[
        x = \frac{f(x)}{f(x_0)}x_0 + \left(x - \frac{f(x)}{f(x_0)}x_0\right).
    \]
    Но
    \[
        x - \frac{f(x)}{f(x_0)x_0} \in \ker f,
    \]
    а $\frac{f(x)}{f(x_0)x_0}$ лежит в линейной оболочке $x_0$.

    Таким образом
    \[
        X = \Lin\{x_0\} \oplus \ker f.
    \]
    и $\forall x \in X \ \exists! y \in \ker f\colon x = y + t x_0$.

    Предположим, что $f\colon X \to \R$ -- линейный непрерывный функционал (эквивалентно, непрерывен в нуле).
    Это действительно очевидно, поскольку если $f$ -- непрерывен в нуле, то есть
    \[
        \forall \epsilon > 0 \exists \delta > 0 \colon \|x\| < \delta \Rightarrow |f(x)| < \epsilon,
    \]
    то для всякой $y \in X\colon$
    \[
        \forall \epsilon > 0 \exists \delta > 0 \colon \|x - y\| < \delta \Rightarrow |f(x - y)| < \epsilon \Longleftrightarrow |f(x) - f(y)| < \epsilon.
    \]


    Аналогично, эквивалентным определением непрерывности $f$ является его ограниченность, то есть $fB_1(0)$ -- ограниченное множество $\Longleftrightarrow$ $f$ -- непрерывный функционал.

    Если $fB_1(0)$ -- ограничено, то $fB_1(0) \subset O_R(0)$
    \[
        \forall \epsilon > 0 \exists \delta = \frac{\epsilon}{R}\colon \|x\| \leq \frac{\epsilon}{R} \Rightarrow |fx| \leq \epsilon.
    \]
    Что и означает непрерывность в 0.

    Если $f$ -- непрерывен, то
    \[
        \forall \epsilon = 1 \exists \delta > 0 \colon fO_{\delta}(0) \subset O_1(0) \Rightarrow fB_{\delta/2}(0) \subset O_1(0).
    \]
\end{note}
\begin{definition}
    Нормой линейного функционала $f$ называется
    \[
        \|f\| = \sup_{\|x\| \leq 1} |f(x)|.
    \]

    По вышепоказанному, $\|f\| < \infty \Longleftrightarrow f$ -- непрерывен.

    Очевидно, что
    \[
        \|f\| = \sup_{\|x\| = 1} |fx| = \sup_{x \neq 0} \frac{|fx|}{\|x\|}.
    \]
\end{definition}
\begin{example}
    Пусть $f$ -- непрерывный линейный функционал, не равный $0$ тождественно.
    Тогда
    \[
        X = \ker f \oplus \Lin\{x_0\},
    \]
    где $x_0 \in X \setminus \ker f$.

    Используя разложение элемента $x \in X\colon$
    \[
        \|f\| = \sup_{t \in \R, y \in \ker f, \ t x_0 + y \neq 0} \frac{|f(tx_0 + y)|}{\|tx_0 + y\|} = \sup_{t > 0, y \in \ker f} \frac{|f(x_0)|}{\frac{1}{t}\|t x_0 + y\|} = \frac{|f(x_0)|}{\inf\limits_{y \in \ker f} \|x_0 + y\|} = \frac{|f(x_0)|}{\rho(x_0, \ker f)}.
    \]
    Поскольку ядро -- замкнута, то $\rho(x_0, \ker f) > 0$.

    Таким образом
    \[
        \rho(x_0, \ker f) = \frac{|f(x_0)|}{\|f\|}.
    \]
\end{example}
\section{Немного о проекциях.}
\begin{definition}
    Говорят, что $\|f\|$ достигается на единичной сфере, если существует $x_0 \in X$ такой, что $\|x_0\| = 1$ и $\|f\| = |f(x_0)|$.
\end{definition}

\begin{note}
    Если $\|f\|$ достигается на сфере, то $x_0$ играет роль <<перпендикуляра>> к $\ker f$.
\end{note}

\begin{example}
    Построим функционал, не имеющий перпендикуляра, то есть такой, что
    \[
        \|f\| > f(x_0) \ \forall x_0\colon  \|x_0\| = 1 \Longleftrightarrow \rho(x_0, \ker f) < 1 \forall x_0\colon \|x_0\| = 1.
    \]

    Пусть $X = C[-1, 1]$ с $\infty$-нормой и определим функционал как
    \[
        f(x) = \int_0^1 x(t) dt - \int_{-1}^0 x(t) dt = \int_{-1}^{1} \sgn(t)x(t)dt.
    \]
    Он, очевидно, ограничен:
    \[
        |f(x)| \leq 2\|x\|.
    \]
    Его ядром являются в точности функции $x \in C[-1, 1]$ такие, что
    \[
        \int_0^{1} x(t) dt = \int_{-1}^0 x(t)dt.
    \]
    Рассмотрим функции вида
    \[
        s_{\epsilon} =
        \begin{cases}
            \sgn(t), & |t| > \epsilon, \\
            \frac{t}{\epsilon}, & |t| \leq \epsilon.
        \end{cases}
    \]
    Такие функции лежат на единичной сфере в пространстве $C[-1, 1]$.
    Функционал $f$ от них равен:
    \[
        2 \geq |f(s_\epsilon)| \geq 2 - \epsilon.
    \]
    Устремляя $\epsilon$ к 0 получаем, что $\|f\| = 2$.

    Пусть существует $x_0 \in C[-1, 1]$ такой, что $\|x_0\| = 1$ и $|f(x_0)| = 2$.
    Без ограничения общности $f(x_0) > 0$.
    То есть
    \[
        \int_{0}^{1} x(t)dt - \int_{-1}^{0} x(t)dt = 2 = \int_{-1}^0 dt + \int_0^1 dt.
    \]
    А значит
    \[
        \int_{-1}^0 1 + x(t)dt + \int_0^1 1 - x(t)dt = 0,
    \]
    и тогда функция должна быть $x(t) \equiv 1$ и $x(t) \equiv -1$ -- чего не может быть.
\end{example}
\begin{example}
    Пусть $f\colon L^1[0, 1] \to \R$ в $(L^1([0, 1]), \|\cdot\|_1)$ определён как
    \[
        fx = \int_0^1 t x(t)dt.
    \]
    Можно показать, что $\|f\| = 1 > |fx|$ при всяком $x \in L^1([0, 1])$ единичной нормы.
\end{example}
\begin{note}[Критерий существования метрической проекции на ядро]
    Заметим, что для непрерывного линейного функционала достижимость своего значения на сфере эквивалентна тому, что для всякого вектора не из ядра существует вектор из ядра, на котором достигается расстояние до ядра.
    То есть, если $f\colon X \to \R$ -- линейный непрерывный функционал, то $\|f\|$ достигается на единичной сфере тогда и только тогда
    \[
        \forall x \in X \setminus \ker f \ \exists z \in \ker f\colon \|x - z\| = \rho(x, \ker f).
    \]
    Что эквивалентно
    \[
        \exists x \in X \setminus \ker f \ \exists z \in \ker f\colon \|x - z\| = \rho(x, \ker f).
    \]
    \begin{proof}
        Пусть $\|f\|$ достигается на единичной сфере, то есть $\exists x_0\colon \|x_0\| = 1 \wedge \|f\| = |f(x_0)|$.
        Возьмём $x \in X \setminus \ker f$.
        Тогда при $t = \frac{-f(x)}{f(x_0)}$ вектор
        \[
           z = x + t x_0 \in \ker f.
        \]
        И
        \[
            \|z - x\| = \left\|\frac{f(x)}{f(x_0)}x_0\right\| = \frac{|f(x)|}{\|f\|} = \rho(x, \ker f).
        \]

        Пусть теперь $\exists x \in X \setminus \ker f \ \exists z \in \ker f \colon \|x - z\| = \rho(x, \ker f)$.
        Рассмотрим
        \[
            x_0 = \frac{x - z}{\|x - z\|}.
        \]
        Тогда
        \[
            |f(x_0)| = \frac{|f(x - z)|}{\|x - z\|} = \frac{|f(x - z)|}{|f(x)|}\|f\| = \|f\|.
        \]
    \end{proof}
\end{note}
\begin{theorem}[Рисс, о проекции]
    Пусть $H$ -- гильбертово пространство и $M \subset H$ -- выпуклое и замкнутое множество в $H$.

    Тогда $\forall x \in H \ \exists ! y \in M\colon \|x - y\| = \rho(x, M)$.
\end{theorem}
\begin{note}
    Множества, для которых в метрических пространствх выполнены условия теоремы Рисса о проекции называются Чёбышовскими.
\end{note}
\begin{proof}
    Заметим, что для всякого $x \in H$ существует минимизирующая последовательность $\{y_n\}$, то есть такая, что
    \[
        \|x - y_n\| \rightarrow \rho(x, M).
    \]

    Согласно критерию Йордана-фон-Неймана для $\|\cdot\|$ выполняется тождество параллелограмма, то есть
    \[
        2\|x\|^2 + 2\|y\|^2 = \|x - y\|^2 + \|x + y\|^2.
    \]
    Тогда
    \[
        \|y_n - y_m\|^2 = \|(y_n - x) - (y_m - x)\|^2 = 2\|y_n - x\|^2 + 2\|y_m - x\|^2 - 4\|\frac{y_n + y_m}{2} - x\|^2.
    \]
    Заметим, что в силу выпуклости $\frac{y_n + y_m}{2} \in M$.
    Тогда $\|\frac{y_n + y_m}{2} - x\|^2 \geq \rho(x, M)^2$.
    Правая часть стремится к 0, при $m,n \rightarrow \infty$ и, таким образом, $\{y_n\}$ -- фундаментальная последовательность.
    Следовательно $y_n \rightarrow y \in M$ -- поскольку $M$ -- замкнуто и $y$ -- искомая точка.

    Единственность проекции показывается аналогичным способом при помощи выпуклости, тождества параллелограмма и предположения о том, что может быть хотя бы две такие точки.

    Пусть $x \in H$, $y$ -- его проекция на $M$ и $z \in M$.
    Тогда
    \[
        \|x - y\|^2 \leq \|x - (y + t(z - y))\|^2, \forall t \in [0, 1].
    \]
    То есть
    \[
        \|x - y\|^2 \leq \|(x - y) + t(y - z)\|^2.
    \]
    Перенесём в правую часть и раскроем скалярные произведения:
    \[
        0 \leq t\langle x - y, y - z \rangle + t^2 \|y - z\|^2.
    \]
    Поделим на $t$ и устремим к $+0\colon$
    \[
        \langle x - y, y - z \rangle \geq 0 \Longleftrightarrow \langle x - y, z - y \rangle \leq 0.
    \]

    Теперь, если у нас есть две точки $x$ и $\tilde{x}$ и, соответсвенно, их проекции $y$ и $\tilde{y}$, то
    \[
        \langle x - y, \tilde{y} - y \rangle \leq 0, \quad \langle \tilde{x} - \tilde{y}, y - \tilde{y} \rangle \leq 0.
    \]
    Сложим эти два неравенства:
    \[
        \langle x - y - \tilde{x} + \tilde{y}, \tilde{y} - y \rangle \leq 0.
    \]
    То есть
    \[
        \langle x - \tilde{x}, \tilde{y} - y \rangle + \|\tilde{y} - y\|^2 \leq 0.
    \]
    Иначе:
    \[
        \|\tilde{y} - y\|^2 \leq \langle \tilde{x} - x, \tilde{y} - y \rangle \leq \|\tilde{x} - x\|\|\tilde{y} - y\|.
    \]
    Что и даёт искомую непрерывность метрической проекции.
\end{proof}
\begin{note}
    Можно заметить, что неравенство $\langle x - y, y - z \rangle \geq 0$ является критерием, то есть если
    \[
        \forall z \in M\colon \|x - z\|^2 = \|x - y + y - z\|^2 = \|x - y \|^2 + 2 \langle x - y, y - z \rangle + \|y - z\|^2 \geq \|x - y\|^2,
    \]
    то точка $y$ является искомой ближайшей точкой.
\end{note}
\begin{corollary}
    Если $L$ -- замкнутое подпространство гильбертового пространства, то $\forall x \in X \setminus L$
    следующие условия эквивалентны:
    \begin{enumerate}
        \item $y$ -- проекция $x$ на $L$;
        \item $\forall z \in L \hookrightarrow \langle x - y, y - z \rangle \geq 0$.
    \end{enumerate}
    Но так как $L$ -- подпространство, то второе условие эквивалентно
    \[
        \langle x - y, \xi \rangle \geq 0,\quad  \forall \xi \in L.
    \]
    С другой стороны, это верно и для $-\xi$ -- а значит условия эквивалентно тому, что
    \[
        \langle x - y, \xi \rangle = 0, \quad \forall \xi \in L.
    \]
    То есть $x - y \in L^{\perp}$.

    Рассмотрим оператор ортогонального проектирования, то есть $P\colon H \to L$ такой, что
    \[
        \langle x - Px, \xi \rangle = 0, \quad \forall \xi \in L.
    \]
    Он, очевидно, линеен.
\end{corollary}
\begin{theorem}[Рисс, об ортогональной проекции]
    Если $L$ -- замкнутое подпространство гильбертового пространства $H$, то $L \oplus L^{\perp} = H$.
\end{theorem}
\begin{note}
    Заметим, что $\|Px\| \leq \|x\|$, так как
    \[
        \|x\|^2 = \|Px\|^2 + \|x - Px\|^2.
    \]
\end{note}